\documentclass[10pt]{article}
\usepackage[a4paper,margin=0.88in]{geometry}
\usepackage{amsmath,amssymb,amsthm,mathtools,enumitem,xcolor,hyperref}
\hypersetup{colorlinks=true,linkcolor=blue!55!black,citecolor=blue!55!black,urlcolor=blue!55!black}
\setlength{\emergencystretch}{3em}
\setlength{\parindent}{0pt}
\setlength{\parskip}{0.55em}
\setlist{leftmargin=1.6em,itemsep=0.25em,topsep=0.25em}
\sloppy

\newtheorem{theorem}{Theorem}[section]
\newtheorem{lemma}[theorem]{Lemma}
\newtheorem{proposition}[theorem]{Proposition}
\newtheorem{corollary}[theorem]{Corollary}
\theoremstyle{definition}
\newtheorem{definition}[theorem]{Definition}
\newtheorem{convention}[theorem]{Convention}
\newtheorem{remark}[theorem]{Remark}

\newcommand{\Fq}{\mathbf F_q}
\newcommand{\Fbar}{\overline{\mathbf F}_q}
\newcommand{\Ql}{\mathbf Q_\ell}
\newcommand{\Qlb}{\overline{\mathbf Q}_\ell}
\newcommand{\Aone}{\mathbf A^1}
\newcommand{\Pone}{\mathbf P^1}
\newcommand{\Spec}{\operatorname{Spec}}
\newcommand{\Gal}{\operatorname{Gal}}
\newcommand{\Tr}{\operatorname{Tr}}
\newcommand{\FT}{\operatorname{FT}}
\newcommand{\RGamma}{R\Gamma}
\newcommand{\RGammac}{R\Gamma_c}
\newcommand{\OO}{\mathcal O}
\newcommand{\Hom}{\operatorname{Hom}}
\newcommand{\Kar}{\operatorname{Kar}}

\begin{document}
\begin{center}
{\large\bfseries The affine-line weight estimate\par}
\medskip
{\normalsize Fourier transform, stationary phase, and cohomological weights\par}
\end{center}

\begin{abstract}
Laumon's affine-line estimate concerns compactly supported cohomology of a
type-3 lisse Weil sheaf.  The coefficient categories, Frobenius and Fourier
normalizations, real weights, middle extensions, and local exact sequence
are specified before the global and local conclusions are stated.
\end{abstract}

\tableofcontents

\section{Arithmetic setting and Frobenius}

Let \(k=\Fq\), where \(q=p^a\), and choose a prime
\(\ell\ne p\).  Fix an algebraic closure \(\Qlb\) of \(\Ql\), a field
embedding
\[
  \iota:\Qlb\longrightarrow\mathbf C,
\]
and a nontrivial additive character
\(\psi:(k,+)\to\Qlb^\times\).  Put
\(E_\psi=\Ql(\psi(k))\), and let \(E/\Ql\) be a finite extension inside
\(\Qlb\) containing \(E_\psi\).  Arithmetic irreducibility below is always
tested after this coefficient field has been chosen.
\begin{definition}[Pro-etale coefficient sheaves]
Let \(\OO_E\) be the ring of integers of \(E\), choose a uniformizer
\(\varpi\), and put \(\Lambda_n=\OO_E/\varpi^n\).  On the small pro-etale
site of a scheme \(X\), let \((\Lambda_n)_X\) be the sheafification of
\[
 U\longmapsto
 \operatorname{Cont}\bigl(\pi_0(|U|),\Lambda_n\bigr).
\]
The integral and rational coefficient sheaves are
\[
 (\OO_E)_X:=\varprojlim_n(\Lambda_n)_X,\qquad
 E_X:=\left(\varprojlim_n(\Lambda_n)_X\right)[1/\varpi].
\]
The inverse limit is taken in sheaves of commutative rings, and the second
formula is ordinary localization of that sheaf.  Thus \(E_X\) is not
introduced as an unqualified constant sheaf with value \(E\).
\end{definition}

\begin{definition}[Ambient derived categories]
Let \(\mathsf{Mod}(E_X)\) be the abelian category of sheaves of
\(E_X\)-modules.  Its ordinary derived category is the localization
\[
 D(X,E):=\mathsf{Cpx}(\mathsf{Mod}(E_X))
        [\{\text{quasi-isomorphisms}\}^{-1}].
\]
A quasi-isomorphism is a chain map inducing isomorphisms on the cohomology
sheaves \(\ker d^i/\operatorname{im}d^{i-1}\).  Shifts and distinguished
triangles are induced by shifted complexes and mapping cones.  Define
\(D(X,\OO_E)\) in the same way from \((\OO_E)_X\)-module sheaves.
\end{definition}

\begin{definition}[Derived-complete integral objects]
For \(L\in D(X,\OO_E)\), put
\[
 L_n=L\otimes^L_{(\OO_E)_X}(\Lambda_n)_X.
\]
Before taking an inverse limit, regard every \(L_n\) as an
\((\OO_E)_X\)-complex by restriction of scalars.  The object \(L\) is
\emph{derived \(\varpi\)-complete} if the canonical map
\[
 L\longrightarrow R\!\varprojlim_n L_n
\]
is an isomorphism in the single category \(D(X,\OO_E)\).

The category \(D^b_c(X,\OO_E)\) is the full subcategory of derived-complete
objects for which every \(L_n\) has constructible cohomology, all reductions
vanish outside one common finite interval, and the canonical transition maps
\[
 L_{n+1}\otimes^L_{(\Lambda_{n+1})_X}(\Lambda_n)_X
          \xrightarrow{\sim}L_n
\]
are isomorphisms.  Constructibility at level \(n\) means locally constant
with finite \(\Lambda_n\)-module stalks on the members of a finite
stratification.
\end{definition}

\begin{definition}[Intrinsic rational constructibility and lattices]
An \(E_X\)-module sheaf is lisse of finite rank if it is locally isomorphic
on the pro-etale site to \(E_X^d\) for some finite \(d\).  It is
constructible if its restriction to each member of a finite disjoint
locally closed stratification is lisse of finite rank.  The category
\(D^b_c(X,E)\) is the full
subcategory of \(D(X,E)\) whose cohomology sheaves are constructible and
vanish outside a finite interval.

An \(\OO_E\)-lattice in \(K\in D^b_c(X,E)\) is an object
\(L\in D^b_c(X,\OO_E)\) together with an isomorphism
\[
          L\otimes^L_{(\OO_E)_X}E_X\xrightarrow{\sim}K.
\]
\end{definition}

\begin{theorem}[Integral-rational comparison]
Assume that \(X\) is noetherian of finite Krull dimension and that
\(\ell\) is invertible on \(X\).  Every object of \(D^b_c(X,E)\) admits an
integral lattice.  For \(L,M\in D^b_c(X,\OO_E)\), rationalization induces
\[
 \Hom_{D^b_c(X,\OO_E)}(L,M)\otimes_{\OO_E}E
 \xrightarrow{\sim}
 \Hom_{D^b_c(X,E)}
   (L\otimes^L_{(\OO_E)_X}E_X,
    M\otimes^L_{(\OO_E)_X}E_X).
\]
Consequently the intrinsic category above is canonically equivalent to
\[
 D^b_c(X,\OO_E)[1/\varpi]^\natural
       \xrightarrow{\sim}D^b_c(X,E),
\]
where the left side first tensors all Hom groups with \(E\) and then takes
the Karoubi completion.  Thus the rational category is exactly the category
obtained from derived-complete integral lattices, including their formal
direct summands.
\end{theorem}

\begin{definition}[Derived operations]
For a morphism \(f:X\to Y\), derived pullback, ordinary derived direct image,
and derived tensor product are the triangulated functors
\[
 Lf^*:D(Y,E)\to D(X,E),\qquad
 Rf_*:D(X,E)\to D(Y,E),\qquad
 -\otimes_{E_X}^L-:D(X,E)\times D(X,E)\to D(X,E),
\]
computed respectively by \(K\)-flat and \(K\)-injective resolutions.  If
\(f\) is separated and of finite type, choose a compactification
\(X\xrightarrow{u}\overline X\xrightarrow{\bar f}Y\), with \(u\) open and
\(\bar f\) proper, and define
\[
 Rf_!=R\bar f_*u_!:D(X,E)\longrightarrow D(Y,E).
\]
Here \(u_!\) is extension by zero.  Compactification independence identifies
the resulting functors for different choices, compatibly with composition.
The constructibility theorems below state when these ambient functors
restrict to \(D^b_c\).
\end{definition}

For a \(k\)-scheme \(X_0\), put
\(X=X_0\times_k\Fbar\).  Put
\[
 \sigma_{\mathrm{ar}}(z)=z^q,\qquad
 \sigma_{\mathrm{geom}}=\sigma_{\mathrm{ar}}^{-1},
\]
and let \(\gamma_X:X\to X\) be the semilinear automorphism induced by
\(\sigma_{\mathrm{geom}}\): \(\gamma_X=1_{X_0}\times\operatorname{Spec}(\sigma_{\mathrm{geom}})\),
so that pulling back a function along \(\gamma_X\) applies \(\sigma_{\mathrm{geom}}\) to its
coefficients in \(\Fbar\).  The \(q\)-power scheme Frobenius
\(\operatorname{Fr}_{X_0,q}:X_0\to X_0\), locally \(f\mapsto f^q\), is a
different morphism and is not denoted by \(\gamma_X\).
The pullback of coefficient modules below uses the canonical identification
\(\gamma_X^{-1}E_X\simeq E_X\).

\begin{definition}[The abelian category of Weil sheaves]
A constructible Weil \(E_X\)-sheaf on \(X_0\) is a pair
\((\mathcal F,\varphi_{\mathcal F})\), where \(\mathcal F\) is a
constructible \(E_X\)-module sheaf on \(X\) and
\[
  \varphi_{\mathcal F}:\gamma_X^*\mathcal F\xrightarrow{\sim}\mathcal F
\]
is an isomorphism.  A morphism \(u:(\mathcal F,\varphi_{\mathcal F})\to
(\mathcal G,\varphi_{\mathcal G})\) is a sheaf morphism satisfying
\[
 u\circ\varphi_{\mathcal F}
   =\varphi_{\mathcal G}\circ\gamma_X^*u.
\]
Kernels and cokernels are formed on the underlying sheaves with their
induced Frobenius isomorphisms.  These objects therefore form an abelian
category \(\mathsf{Cons}_{\mathrm W}(X_0,E)\).  A Weil sheaf is lisse when
its underlying \(E_X\)-module sheaf is lisse of finite rank.
A lisse Weil sheaf is \emph{arithmetically irreducible} if it has no nonzero
proper Weil subsheaf.  This does not imply that its geometric pullback is
irreducible.
\end{definition}

\begin{definition}[Strict Weil complexes and their localization]
A strict bounded constructible Weil complex is a bounded cochain complex in
\(\mathsf{Cons}_{\mathrm W}(X_0,E)\); equivalently, every differential is a
morphism of Weil sheaves.  Equivariant chain homotopies define the homotopy
category
\[
 \mathsf K^b(\mathsf{Cons}_{\mathrm W}(X_0,E)).
\]
Equivariant shifts and mapping cones make it an ordinary triangulated
category.  A chain map is an equivariant quasi-isomorphism if its underlying
map of \(E_X\)-module complexes is a quasi-isomorphism.  The finite-stage
Weil derived category is the Verdier localization
\[
 \operatorname{WeilD}^b_c(X_0,E)
 :=\mathsf K^b(\mathsf{Cons}_{\mathrm W}(X_0,E))
       [\{\text{equivariant quasi-isomorphisms}\}^{-1}].
\]
Its shift and distinguished triangles are induced by equivariant shifts
and mapping cones.  Accordingly, an object is represented by a strict
equivariant complex modulo equivariant quasi-isomorphism.  A derived
isomorphism \(\gamma_X^*K\simeq K\) on an already-localized object does not,
without a strict equivariant representative, define an object of this
category.
\end{definition}

\begin{theorem}[Forgetting and lifting Frobenius]
Forgetting the Weil structure gives an exact conservative functor
\[
 \operatorname{For}_E:\operatorname{WeilD}^b_c(X_0,E)
             \longrightarrow D^b_c(X,E).
\]
Every source object has an induced derived isomorphism
\(\gamma_X^*\operatorname{For}_E(K)\simeq\operatorname{For}_E(K)\).
An exact operation represented on complexes and equipped with a chain-level
comparison with \(\gamma_X^*\) descends to the localization and hence
defines an exact operation on Weil complexes.  In particular, this applies
to coefficient-compatible six operations and Fourier transforms.
\end{theorem}

\begin{definition}[Iterated Weil datum and rational-point Frobenius]
For a Weil sheaf \((\mathcal F,\varphi_{\mathcal F})\) and \(n\geq1\), put
\[
 \varphi_{\mathcal F}^{[n]}:=
 \varphi_{\mathcal F}\circ\gamma_X^*(\varphi_{\mathcal F})\circ
 \cdots\circ\gamma_X^{(n-1)*}(\varphi_{\mathcal F}):
 \gamma_X^{n*}\mathcal F\xrightarrow{\sim}\mathcal F.
\]
For \(x\in X_0(\mathbf F_{q^n})\), choose a geometric point \(\bar x\)
above \(x\).  The \(\mathbf F_{q^n}\)-rational structure identifies
\(\gamma_X^n(\bar x)\) with \(\bar x\); the stalk of
\(\varphi_{\mathcal F}^{[n]}\) is the geometric Frobenius operator
\[
 F_{x,n}(\mathcal F)\in
 \operatorname{Aut}_E(\mathcal F_{\bar x}).
\]
If \(|x|\in|X_0|\) is the closed point below \(x\), of degree
\(d_x=[\kappa(|x|):k]\mid n\), define
\(F_{|x|}:=F_{x,d_x}\).  Up to the canonical conjugacy arising from the
choice of \(\bar x\), one has
\[
                   F_{x,n}=F_{|x|}^{n/d_x}.
\]
For a strict Weil complex \(K\), the same construction on each stalk
cohomology group defines \(F_{x,n}^i(K)\), and
\[
 t_K^{(n)}(x):=\sum_i(-1)^i
     \operatorname{Tr}\bigl(F_{x,n}^i(K)\mid\mathcal H^i(K)_{\bar x}\bigr)
\]
defines a function \(t_K^{(n)}:X_0(\mathbf F_{q^n})\to E\).
For compactly supported cohomology, geometric Frobenius is induced by the
same iterated Weil datum and compact-support base change; it is not obtained
by applying cohomology to \(\operatorname{Fr}_{X_0,q}\).
\end{definition}

\section{Common finite stages and \texorpdfstring{$\Qlb$}{Qbar-l}}

Let \(\mathsf{FinExt}_\ell\) be the filtered poset of finite extensions
\(E\subset\Qlb\).  For \(E\subset M\), scalar extension is the exact functor
\[
 \operatorname{ext}_{M/E}:D^b_c(X,E)\longrightarrow D^b_c(X,M),
 \qquad K\longmapsto K\otimes^L_{E_X}M_X.
\]
Derived tensor associativity supplies the coherent comparisons for towers
\(E\subset M\subset N\).

\begin{definition}[Common-stage Hom groups]
For \(E,E'\in\mathsf{FinExt}_\ell\), let \(\mathsf I(E,E')\) be the filtered
poset of finite extensions \(M/\Ql\) inside \(\Qlb\) that contain both.  The finite-stage
pair category \(\mathcal C_X^{\mathrm{pair}}\) has objects \((E,K)\) with
\(K\in D^b_c(X,E)\) and
\[
\begin{aligned}
 \Hom_{\mathcal C_X^{\mathrm{pair}}}((E,K),(E',L))
 :=\operatorname*{colim}_{M\in\mathsf I(E,E')}
 \Hom_{D^b_c(X,M)}
 \bigl(\operatorname{ext}_{M/E}K,\operatorname{ext}_{M/E'}L\bigr).
\end{aligned}
\]
A morphism is represented over one common finite coefficient extension.
Two representatives agree when their scalar extensions agree over a
further common finite coefficient extension.  Addition and composition are performed after moving
all representatives involved to one common finite stage.
\end{definition}

\begin{definition}[Karoubi triples]
For an additive category \(\mathcal C\), its Karoubi envelope
\(\Kar(\mathcal C)\) has objects \((P,e)\) with \(e^2=e\).  A morphism
\(u:(P,e)\to(Q,d)\) satisfies \(u=d\,u\,e\), and the identity of
\((P,e)\) is \(e\).  Define
\[
 D^b_c(X,\Qlb)_{\mathrm{2colim}}
      :=\Kar(\mathcal C_X^{\mathrm{pair}}).
\]
Thus an object is a triple \((E,K,e)\), where \(K\in D^b_c(X,E)\) and
\(e\) is an idempotent represented at a finite common stage.  Its shift is
\((E,K[1],e[1])\), and a triangle is distinguished when it is isomorphic to
a direct summand of a triangle that is distinguished at one common finite
stage.
\end{definition}

\begin{definition}[The \(\Qlb\)-coefficient Weil category]
Scalar extension preserves strict Weil complexes and equivariant
quasi-isomorphisms.  Form the Weil pair category by replacing
\(D^b_c(X,E)\) in the common-stage Hom formula with
\(\operatorname{WeilD}^b_c(X_0,E)\), and requiring every representative to
commute with Frobenius.  Explicitly,
\[
\begin{aligned}
 &\Hom_{\mathcal C_{X_0}^{\mathrm W,pair}}((E,K),(E',L))\\
 &\quad:=
 \operatorname*{colim}_{M\in\mathsf I(E,E')}
 \Hom_{\operatorname{WeilD}^b_c(X_0,M)}
 \bigl(\operatorname{ext}_{M/E}K,
       \operatorname{ext}_{M/E'}L\bigr).
\end{aligned}
\]
Define
\[
 \operatorname{WeilD}^b_c(X_0,\Qlb)_{\mathrm{2colim}}
 :=\Kar(\mathcal C_{X_0}^{\mathrm W,pair}).
\]
Its objects are triples \((E,K,e)\) with
\(K\in\operatorname{WeilD}^b_c(X_0,E)\), its morphisms are
common-stage equivariant morphisms \(u\) satisfying \(u=d\,u\,e\), and its
triangles are the formal retracts of stagewise distinguished triangles.
\end{definition}

\begin{theorem}[Triangulation and coefficient independence]
The two pair categories are ordinary triangulated categories under their
stagewise shifts and triangles.  Their Karoubi envelopes are
idempotent-complete ordinary triangulated categories.  Every finite-stage
scalar-extension functor is exact, and any equation between finitely many
objects and morphisms can be checked after extension to one common finite
extension of \(\Ql\).
\end{theorem}

\begin{theorem}[Cohomology and descent of operations]
For a Karoubi triple \((E,K,e)\), its \(i\)-th cohomology is the
image of the idempotent \(H^i(e)\) on the finite-stage cohomology sheaf
\(H^i(K)\).  This definition is unchanged after passage to a larger common
finite extension of \(\Ql\).  A family of finite-stage exact functors equipped with coherent
scalar-extension comparisons descends to the pair categories and their
Karoubi completions.  The same holds equivariantly for functors represented
on strict Weil complexes.  Thus compactly supported cohomology, middle
extension, and Fourier transform define exact functors on the corresponding
\(\Qlb\)-coefficient categories.
\end{theorem}

\section{Real weights}

\begin{definition}[Pure eigenvalue and upper weight bound]
Let \(V\) be a finite-dimensional \(E\)-vector space with an automorphism
\(F\), and let \(Q>1\).  For a real number \(w\), an eigenvalue
\(\alpha\in\Qlb\) of \(F\) is \emph{\(\iota\)-pure of weight \(w\)
over \(Q\)} if
\[
  |\iota(\alpha)|=Q^{w/2}.
\]
The Frobenius module \(V\) is \(\iota\)-pure of weight \(w\) if every
eigenvalue has this absolute value.  It has \(\iota\)-weights at most \(v\)
if every eigenvalue satisfies
\[
  |\iota(\alpha)|\le Q^{v/2}.
\]
It has \(\iota\)-weights at least \(v\) if every eigenvalue satisfies
\[
  |\iota(\alpha)|\ge Q^{v/2}.
\]
The numbers \(w\) and \(v\) are real numbers; no integrality assumption is
built into these definitions.
\end{definition}

\begin{definition}[Punctual purity and mixedness]
A constructible Weil sheaf \(\mathcal F\) on \(X_0\) is
\emph{punctually \(\iota\)-pure of weight \(w\in\mathbf R\)} if, for every
closed point \(x\), every eigenvalue \(\alpha\) on the nonzero stalk
\(\mathcal F_{\bar x}\) satisfies
\[
  |\iota(\alpha)|=q_x^{w/2}.
\]
A constructible Weil sheaf is \(\iota\)-mixed of weights at most \(v\) if it
has a finite filtration by Weil subsheaves whose nonzero successive
quotients are punctually \(\iota\)-pure of real weights at most \(v\).
For a bounded constructible Weil complex \(K\), the assertion that \(K\) has
weights at most \(v\) means that \(\mathcal H^m(K)\) is mixed of weights at
most \(v+m\) for every \(m\).
\end{definition}

A weight assertion is relative to the fixed embedding \(\iota\).  The
all-embeddings version requires the corresponding assertion
for every embedding \(\Qlb\to\mathbf C\); it is not a consequence of
a single fixed-\(\iota\) inequality.

\begin{convention}[Tate twists]
Geometric Frobenius acts on \(E(1)\) by \(q^{-1}\).  Thus a Tate twist
\((n)\) lowers weights by \(2n\), and \((-1)\) raises weights by \(2\).
\end{convention}

\section{Ordinary and perverse middle extension}

Put \(A_0=\Aone_k\), \(A=\Aone_{\Fbar}\), and let
\(j_0:U_0\hookrightarrow A_0\) be a dense open immersion.  Its geometric
base change is denoted \(j:U\hookrightarrow A\).

\begin{definition}[Middle perverse category on a curve]
Let \(C\) be a smooth curve over an algebraically closed field.  For
\(K\in D^b_c(C,E)\), set
\[
 \begin{aligned}
 {}^pD^{\le0}(C,E)
   &=\{K:\dim\operatorname{Supp}\mathcal H^m(K)\le -m
          \text{ for every }m\},\\
 {}^pD^{\ge0}(C,E)
   &=\{K:D_CK\in{}^pD^{\le0}(C,E)\},
 \end{aligned}
\]
where
\[
 D_CK=R\mathcal Hom_C(K,E_C(1)[2])
\]
is Verdier duality, \(E_C\) is the constant \(E\)-sheaf, and the equality
\(E_C(1)[2]\simeq\omega_C\) is the smooth purity identification of the
dualizing complex of \(C\).  The heart
\[
  \operatorname{Perv}(C,E)
  :={}^pD^{\le0}(C,E)\cap{}^pD^{\ge0}(C,E)
\]
is the middle perverse category.  With this normalization, a lisse sheaf
\(\mathcal L\) on a smooth curve is placed in perverse degree zero as
\(\mathcal L[1]\).
\end{definition}

\begin{definition}[Perverse middle extension]
For \(P\in\operatorname{Perv}(U,E)\), its middle extension
\(j_{!*}P\) is the image, in the abelian category
\(\operatorname{Perv}(A,E)\), of the canonical morphism
\[
  {}^p\!\mathcal H^0(Rj_!P)\longrightarrow
  {}^p\!\mathcal H^0(Rj_*P).
\]
Equivalently, it is the unique perverse extension of \(P\) having no
nonzero perverse subobject and no nonzero perverse quotient supported on
\(A\setminus U\).
\end{definition}

\begin{proposition}[Middle extension of a lisse sheaf]
For a lisse \(E\)-sheaf \(\mathcal F\) on \(U\), there is a canonical
isomorphism
\[
  j_{!*}(\mathcal F[1])\xrightarrow{\sim}(j_*\mathcal F)[1].
\]
Here \(j_*\mathcal F=R^0j_*\mathcal F\) is the ordinary sheaf-theoretic
direct image, not the derived object \(Rj_*\mathcal F\).  If
\(\bar s\) is a geometric point of \(A\setminus U\), then
\[
  (j_*\mathcal F)_{\bar s}
  \cong \mathcal F_{\bar\eta_s}^{I_s},
\]
where \(I_s\) is the inertia group of the punctured strict henselian trait at
\(s\).  In particular, \(j_*\mathcal F\) is constructible.
\end{proposition}

\begin{definition}[Weil perverse sheaves and middle extension]
Define
\[
 \begin{aligned}
 {}^pD_{\mathrm W}^{\leq0}(A_0,E)
  &=\{K\in\operatorname{WeilD}^b_c(A_0,E):
        \operatorname{For}_E(K)\in{}^pD^{\leq0}(A,E)\},\\
 {}^pD_{\mathrm W}^{\geq0}(A_0,E)
  &=\{K\in\operatorname{WeilD}^b_c(A_0,E):
        \operatorname{For}_E(K)\in{}^pD^{\geq0}(A,E)\}.
 \end{aligned}
\]
These subcategories form a perverse \(t\)-structure.  Its abelian heart is
denoted \(\operatorname{Perv}_{\mathrm W}(A_0,E)\), and its degree-zero
cohomology functor is \({}^p\mathcal H^0_{\mathrm W}\).  For a dense open
immersion \(j_0:U_0\hookrightarrow A_0\) and
\(P\in\operatorname{Perv}_{\mathrm W}(U_0,E)\), define
\[
 j^W_{0!*}P:=\operatorname{im}_{\operatorname{Perv}_{\mathrm W}(A_0,E)}
 \left({}^p\mathcal H^0_{\mathrm W}(Rj_{0!}P)
       \longrightarrow{}^p\mathcal H^0_{\mathrm W}(Rj_{0*}P)\right).
\]
The forgetful functor is perverse \(t\)-exact and satisfies
\[
 \operatorname{For}_E(j^W_{0!*}P)
       \simeq j_{!*}\operatorname{For}_E(P).
\]
For a lisse Weil sheaf \(\mathcal F_W\), the notation
\(j_{0*}\mathcal F_W\) means the ordinary degree-zero direct-image Weil
sheaf, whose underlying geometric sheaf is
\(j_*\mathcal F=R^0j_*\mathcal F\).  There is a canonical isomorphism
\[
 j^W_{0!*}(\mathcal F_W[1])\simeq(j_{0*}\mathcal F_W)[1].
\]
Thus neither \(j_{0*}\mathcal F_W\) nor \(j_*\mathcal F\) denotes a
derived direct image.
\end{definition}

Let \(a:A\to\Spec\Fbar\) be the structure morphism.  Its ambient exceptional
direct image is the triangulated functor
\[
  Ra_!:D(A,E)\longrightarrow D(\Spec\Fbar,E).
\]

\begin{proposition}[Constructibility and compact-support finiteness]
For every \(K\in D^b_c(A,E)\), one has
\[
  Ra_!K\in D^b_c(\Spec\Fbar,E).
\]
Equivalently, \(Ra_!K\) is a bounded complex of finite-dimensional
\(E\)-vector spaces.  For a finite-rank lisse sheaf \(\mathcal F\) on \(U\),
both \(j_*\mathcal F\) and \(R^1j_*\mathcal F\) are constructible, so this
statement applies in particular to \(j_*\mathcal F\).
\end{proposition}

For a constructible sheaf \(\mathcal G\) on \(A\), define
\[
  \RGammac(A,\mathcal G):=Ra_!\mathcal G\in D^b_c(\Spec\Fbar,E),
  \qquad
  H_c^i(A,\mathcal G):=\mathcal H^i(Ra_!\mathcal G).
\]
When \(\mathcal G\) carries a Weil structure, compact-support base change
for the Frobenius square induces geometric Frobenius on every
\(H_c^i(A,\mathcal G)\).

\section{The type-3 condition}

Let \(\Pone_k=A_0\cup\{\infty\}\).  Write
\(T_\infty^h=\Spec\mathcal O^h_{\Pone_k,\infty}\), let
\(\eta_\infty\) be its generic point, and choose a geometric generic point
\(\bar\eta_\infty\).  Define
\[
  D_\infty:=\pi_1(\eta_\infty,\bar\eta_\infty),
  \qquad
  I_\infty:=\ker\bigl(D_\infty\longrightarrow\Gal(\Fbar/k)\bigr).
\]
Let \(F_q\in\Gal(\Fbar/k)\) be geometric Frobenius and put
\[
 W(k):=\langle F_q\rangle\cong\mathbf Z,
 \qquad
 W_\infty:=D_\infty\times_{\Gal(\Fbar/k)}W(k).
\]
Then
\[
 1\longrightarrow I_\infty\longrightarrow W_\infty
   \longrightarrow\langle F_q\rangle\longrightarrow1.
\]
The topology on \(W_\infty\) makes \(I_\infty\) open with its profinite
topology and gives the quotient the discrete topology.  A lift
\(\Phi_\infty\) of \(F_q\) gives a noncanonical identification
\(W_\infty\cong I_\infty\rtimes\langle\Phi_\infty\rangle\).  The local
monodromy of a lisse Weil sheaf at infinity is the resulting continuous
\(E\)-representation on \(\mathcal F_{\bar\eta_\infty}\); no continuous
action of the whole profinite group \(D_\infty\) is part of a Weil datum.

Fix the standard chart \(\iota:A_0\hookrightarrow\Pone_k\),
\(x\mapsto[x:1]\), whose image is the standard open \(D_+(T_1)\) and whose
complement is the point \(\infty=[1:0]\).

\begin{definition}[Unramified and geometrically constant]
The sheaf \(\mathcal F\) is \emph{unramified at infinity} if it extends
to a lisse Weil sheaf on \(U_0\cup\{\infty\}\): there are an open
\(W\subseteq\Pone_k\) with \(\infty\in W\) and \(\iota^{-1}(W)=j_0(U_0)\), and
a lisse Weil sheaf \(\mathcal G\) on \(W\) with
\(j_0'^*\mathcal G\simeq\mathcal F\), where \(j_0':U_0\hookrightarrow W\)
is the factorisation of \(\iota\circ j_0\) through \(W\).  It is
\emph{geometrically constant} if, on each connected component of the
geometric fiber \(U=U_0\otimes\Fbar\), it is isomorphic to the constant
lisse sheaf associated with some finite-dimensional \(E\)-vector space.
\end{definition}

\begin{remark}
Unramifiedness at infinity does not depend on the chart: two open
immersions \(A_0\hookrightarrow\Pone_k\) over \(k\) differ by an
automorphism of \(\Pone_k\), since \(\Pone_k\) is the smooth
compactification of \(A_0\).  It is equivalent to \(I_\infty\) acting
trivially on \(\mathcal F_{\bar\eta_\infty}\), since a local system on
\(\eta_\infty\) extends to the henselian trait \(T_\infty^h\) exactly when it
is unramified.  Since \(U\) is connected, geometric constancy is equivalent
to the geometric fundamental group acting trivially on a stalk.
\end{remark}

\begin{definition}[Geometric Artin--Schreier subquotients]
For a finite extension \(k'/k\), put
\[
 \psi_{k'}=\psi\circ\Tr_{k'/k}:(k',+)\longrightarrow\Qlb^\times.
\]
A simple subquotient of the geometric local system
\(\mathcal F\otimes_E\Qlb\) is called \emph{Artin--Schreier} if, after
descent to some finite \(k'/k\) and some finite coefficient extension, it
is isomorphic to
\[
 \left.\mathcal L_{\psi_{k'}}(a x)\right|_{U_{\Fbar}}
 \qquad\text{for some }a\in k'.
\]
Here \(\mathcal L_{\psi_{k'}}(a x)\) is the pullback of the
Artin--Schreier sheaf by \(x\mapsto ax\).
\end{definition}

\begin{proposition}[Artin--Schreier alternatives at infinity]
For \(a\ne0\), the geometric sheaf
\(\mathcal L_{\psi_{k'}}(a x)\) has Swan conductor one at infinity.  For
\(a=0\), it is the constant rank-one sheaf.  Consequently, an unramified
geometric local system has no Artin--Schreier simple subquotient if and only
if it has no geometrically constant simple subquotient.
\end{proposition}

\begin{definition}[Type-3 affine-line datum]
Fix \(w\in\mathbf R\).  A lisse Weil sheaf
\(\mathcal F_W=(\mathcal F,\varphi_{\mathcal F})\) on \(U_0\) is called a
\emph{type-3 datum of weight \(w\)} if all of the following hold:
\begin{enumerate}
  \item \(\mathcal F_W\) has positive finite rank and is irreducible in the
  arithmetic Weil category;
  \item \(\mathcal F_W\) is punctually \(\iota\)-pure of weight \(w\);
  \item \(\mathcal F\) is unramified at infinity;
  \item \(\mathcal F\otimes_E\Qlb\) has no geometric
  Artin--Schreier simple subquotient.
\end{enumerate}
Under the unramifiedness condition, the last clause is equivalently the
absence of a geometrically constant simple subquotient.  Geometric
irreducibility is not part of the definition.
\end{definition}

\section{Fourier and local normalizations}

Let \(A'_0\) be the affine line dual to \(A_0\), with coordinate \(x'\) and
origin \(0'\), and put \(A'=A'_0\times_k\Fbar\).  Let
\(\mathcal L_{\psi,E}\) be the rank-one Artin--Schreier Weil \(E\)-sheaf
characterized by
\[
  \Tr\bigl(F_u\mid(\mathcal L_{\psi,E})_{\bar u}\bigr)
  =\psi\bigl(\Tr_{\kappa(u)/k}(u)\bigr)
\]
for every closed point \(u\) of \(\Aone_k\).
For every finite extension \(M/E\) inside \(\Qlb\), put
\(\mathcal L_{\psi,M}=\mathcal L_{\psi,E}\otimes_E M\).  These sheaves are
compatible with further scalar extension.

Let \(p:A\times A'\to A\) and \(p':A\times A'\to A'\) be the projections,
and let \(\mu(x,x')=xx'\).  On the ambient triangulated categories define
\[
  \FT_\psi^{\mathrm{amb}}(K)
  :=Rp'_!\bigl(Lp^*K
       \otimes^L_{E_{A\times A'}}L\mu^*\mathcal L_{\psi,E}\bigr)[1],
  \qquad
  D(A,E)\longrightarrow D(A',E).
\]

\begin{proposition}[Fourier finiteness]
If \(K\in D^b_c(A,E)\), then
\(\FT_\psi^{\mathrm{amb}}(K)\in D^b_c(A',E)\).  Hence the displayed formula
restricts to a triangulated functor
\[
  \FT_\psi:D^b_c(A,E)\longrightarrow D^b_c(A',E).
\]
The shift \([1]\) is part of the definition, and there is no Tate twist in
this normalization.
\end{proposition}

\begin{definition}[Fourier transform of Weil complexes]
Let \(p_0:A_0\times_kA'_0\to A_0\) and
\(p'_0:A_0\times_kA'_0\to A'_0\) be the projections, and let
\(\mu_0(x,x')=xx'\).  For every finite extension \(M/E\) inside
\(\Qlb\), define the exact triangulated functor
\[
 \begin{aligned}
 \FT_{\psi,M,W}:\operatorname{WeilD}^b_c(A_0,M)
   &\longrightarrow\operatorname{WeilD}^b_c(A'_0,M),\\
 K&\longmapsto Rp'_{0!}\left(
     Lp_0^*K\otimes^L_{M_{A_0\times A'_0}}
     L\mu_0^*\mathcal L_{\psi,M}\right)[1].
 \end{aligned}
\]
For \(M=E\), write \(\FT_{\psi,W}\).  The common-stage
\(\Qlb\)-coefficient functor is defined on the cofinal system of finite
coefficient extensions containing \(E_\psi\), by the common-stage Hom
colimit and Karoubi completion.
\end{definition}

\begin{theorem}[Fourier comparison and perverse exactness]
The Weil Fourier transforms commute with finite scalar extension, and
\[
 \operatorname{For}_M\circ\FT_{\psi,M,W}
   \xrightarrow{\sim}
 \FT_{\psi,M}\circ\operatorname{For}_M.
\]
They are perverse \(t\)-exact and induce equivalences
\[
 \FT_{\psi,M,W}:\operatorname{Perv}_{\mathrm W}(A_0,M)
   \xrightarrow{\sim}\operatorname{Perv}_{\mathrm W}(A'_0,M).
\]
If \(\FT^{\mathrm{rev}}_{\psi^{-1},M,W}\) denotes the transform from the
dual line back to the original line with the inverse character, then
\[
 \FT^{\mathrm{rev}}_{\psi^{-1},M,W}\circ\FT_{\psi,M,W}
   \xrightarrow{\sim}(-1),\qquad
 \FT_{\psi,M,W}\circ\FT^{\mathrm{rev}}_{\psi^{-1},M,W}
   \xrightarrow{\sim}(-1).
\]
Thus \(\FT^{\mathrm{rev}}_{\psi^{-1},M,W}(1)\) is a quasi-inverse.
The same assertions hold in the common-stage Karoubi categories.
\end{theorem}

\begin{theorem}[Fourier naturality, coordinates, and finite ground fields]
Let \(m\geq1\), put \(k_m=\mathbf F_{q^m}\), and define
\[
 \psi_m:=\psi\circ\operatorname{Tr}_{k_m/k}:(k_m,+)\longrightarrow E^\times.
\]
On strict Weil complexes define
\[
 \operatorname{Res}_{k_m/k}(K,\phi_K):=(K,\phi_K^{[m]}),
\]
viewed as an object on \(A_{0,m}=A_0\times_k k_m\).  There is a natural
isomorphism of exact triangulated functors
\[
 \operatorname{Res}_{k_m/k}\circ\FT_{\psi,M,W}
 \xrightarrow{\sim}
 \FT_{\psi_m,M,W}\circ\operatorname{Res}_{k_m/k}.
\]
It is compatible with towers \(k_r/k_m/k\), finite scalar extension
\(M\subset M'\), the forgetful functor, perverse truncation, and Fourier
inversion.

For \(a\in k^\times\), let \(m_a:A_0\to A_0\) be multiplication by \(a\)
and let \(m'_{a^{-1}}:A'_0\to A'_0\) be multiplication by \(a^{-1}\).
The equality
\[
 \mu_0\circ(m_a\times m'_{a^{-1}})=\mu_0
\]
induces a natural isomorphism
\[
 (m'_{a^{-1}})^*\FT_{\psi,M,W}(K)
 \xrightarrow{\sim}
 \FT_{\psi,M,W}(m_a^*K).
\]
These isomorphisms satisfy the identity and composition equations for
\(a,b\in k^\times\).  Thus the transform and its later local comparison
maps are independent of a simultaneous change \(x_1=ax\),
\(x'_1=a^{-1}x'\) of dual coordinates.
\end{theorem}

\begin{theorem}[Middle extension of the type-3 input]
Let \(\mathcal F_W\) be a type-3 datum and let
\(j_0:U_0\hookrightarrow A_0\) be its domain.  The ordinary direct image
\(j_{0*}\mathcal F_W\) is a constructible Weil sheaf, and there are
canonical isomorphisms
\[
 j^W_{0!*}(\mathcal F_W[1])
 \xrightarrow{\sim}(j_{0*}\mathcal F_W)[1]
 \quad\text{in }\operatorname{Perv}_{\mathrm W}(A_0,E),
\]
and
\[
 j_{!*}(\mathcal F[1])
 \xrightarrow{\sim}(j_*\mathcal F)[1]
 \quad\text{in }\operatorname{Perv}(A,E).
\]
The second is the image under the forgetful functor of the first.  Both
isomorphisms commute with finite coefficient extension.
\end{theorem}

\begin{proposition}[Cohomological concentration before the estimate]
For a type-3 datum, the absence of a geometrically constant simple
subquotient gives
\[
 H_c^0(A,j_*\mathcal F)=0,
 \qquad H_c^2(A,j_*\mathcal F)=0.
\]
Moreover \(R\Gamma_c(A,j_*\mathcal F)\) is bounded with
finite-dimensional cohomology.  Thus its only possibly nonzero
cohomology group is the finite-dimensional Frobenius module
\(H_c^1(A,j_*\mathcal F)\).
\end{proposition}

For a type-3 sheaf set
\[
 \begin{aligned}
 K_{\mathcal F,W}&:=j^W_{0!*}(\mathcal F_W[1])
              =(j_{0*}\mathcal F_W)[1]
              \in\operatorname{Perv}_{\mathrm W}(A_0,E),\\
 K_{\mathcal F}&:=\operatorname{For}_E(K_{\mathcal F,W})
              =(j_*\mathcal F)[1]
              =j_{!*}(\mathcal F[1])\in\operatorname{Perv}(A,E).
 \end{aligned}
\]
Let \(j':U'=A'\setminus\{0'\}\hookrightarrow A'\).  At \(0'\), define
the arithmetic henselian trait and its generic point by
\[
 T_{0'}^h=\Spec\mathcal O^h_{A'_0,0'},
 \qquad
 \eta_{0'}^h=\Spec\operatorname{Frac}(\mathcal O^h_{A'_0,0'}).
\]
Choose a geometric generic point \(\bar\eta_{0'}\) and define
\[
  D_{0'}:=\pi_1(\eta_{0'}^h,\bar\eta_{0'}),
  \qquad
  I_{0'}:=\ker\bigl(D_{0'}\longrightarrow\Gal(\Fbar/k)\bigr).
\]
Define
\[
 W_{0'}:=D_{0'}\times_{\Gal(\Fbar/k)}W(k).
\]
It fits into
\[
 1\longrightarrow I_{0'}\longrightarrow W_{0'}
   \longrightarrow\langle F_q\rangle\longrightarrow1.
\]
Its topology is defined as for \(W_\infty\), and a lift \(\Phi_{0'}\) of
\(F_q\) gives a noncanonical identification
\(W_{0'}\cong I_{0'}\rtimes\langle\Phi_{0'}\rangle\).  Since \(0'\) and
\(\infty\) are \(k\)-rational, both quotient groups
\(W_{0'}/I_{0'}\) and \(W_\infty/I_\infty\) are canonically \(W(k)\).

\begin{theorem}[Type-3 Fourier classification]
Let \(\mathcal F_W\) be a type-3 datum.  There is an arithmetically
irreducible lisse Weil \(E\)-sheaf \(\mathcal F'_W\) on
\(U'_0=A'_0\setminus\{0'\}\), with geometric sheaf \(\mathcal F'\), and a
canonical isomorphism in the Weil perverse category
\[
  \FT_{\psi,W}(K_{\mathcal F,W})
  \xrightarrow{\sim}j'^W_{0!*}(\mathcal F'_W[1])
  \xrightarrow{\sim}(j'_{0*}\mathcal F'_W)[1].
\]
After applying \(\operatorname{For}_E\), this is the geometric isomorphism
\[
 \FT_\psi(K_{\mathcal F})
  \xrightarrow{\sim}j'_{!*}(\mathcal F'[1])
  \xrightarrow{\sim}(j'_*\mathcal F')[1].
\]
The transformed lisse sheaf is recovered from the transform by
\[
 \mathcal F'
 =\mathcal H^{-1}\!\left(
   j'^*\operatorname{For}_E\FT_{\psi,W}(K_{\mathcal F,W})\right).
\]
The transformed perverse sheaf has no nonzero direct summand supported at a
closed point.  No geometric irreducibility assertion for \(\mathcal F'\) is
included.
\end{theorem}

\begin{definition}[Purity of the transformed lisse sheaf]
Let \(\mathcal F'_W\) be the arithmetically irreducible lisse Weil sheaf
defined by the type-3 Fourier-classification theorem.  It is
\emph{punctually \(\iota\)-pure of weight \(w'\in\mathbf R\)} if, for every
closed point \(x'\in|U'_0|\), every eigenvalue \(\beta\) of geometric
Frobenius on \(\mathcal F'_{\bar x'}\) satisfies
\[
                   |\iota(\beta)|=q_{x'}^{w'/2}.
\]
This property concerns \(\mathcal F'_W\) alone; it does not include the
choice or existence of a companion for \(\mathcal F_W\).
\end{definition}

\begin{theorem}[Unramified boundary module]
Let
\[
 V_\infty=\mathcal F_{\bar\eta_\infty}
 \in\operatorname{Rep}_E(W_\infty).
\]
Then \(I_\infty\) acts trivially on \(V_\infty\), so
\(V_\infty=(V_\infty)^{I_\infty}\) and the action factors through
\(W_\infty/I_\infty\simeq W(k)\).  The module \(V_\infty\) is nonzero and
is \(\iota\)-pure of weight \(w\).  With geometric-Frobenius Tate
normalization, \(V_\infty(-1)\) is nonzero and pure of weight \(w+2\).
After identifying \(W_\infty/I_\infty\) with
\(W_{0'}/I_{0'}\), it is an unramified \(W_{0'}\)-representation.
\end{theorem}

\begin{definition}[Affine-boundary-real companion]
Let \(\mathcal F_W^c\) be a type-3 datum of weight \(w\) on \(U_0\).  It is a
\emph{boundary-real \(\iota\)-companion on \(\mathbf A^1\)} of
\(\mathcal F_W\) if the following two conditions hold.
\begin{enumerate}
\item At every closed point \(u\in|U_0|\), the Frobenius eigenvalue multiset
      of \(\mathcal F_W^c\), after applying \(\iota\), is the complex
      conjugate of that of \(\mathcal F_W\).
\item For every finite boundary point
      \(s\in\mathbf A^1_k\setminus U_0\), with inertia group
      \(I_s\) and geometric Frobenius \(F_s\),
      \[
       \iota\det\!\left(1-TF_s\mid
          (\mathcal F^c_{\bar\eta_s})^{I_s}\right)
       =\overline{\iota\det\!\left(1-TF_s\mid
          (\mathcal F_{\bar\eta_s})^{I_s}\right)}.
      \]
\end{enumerate}
The second condition is part of the definition and is stronger than the
existence of a conjugate companion on the open curve.  No reality condition
at \(\infty\) is imposed; condition (2) concerns exactly the finite boundary
stalks of \(j_{0*}\) on \(\mathbf A^1_k\).
\end{definition}

\begin{theorem}[Conjugate-dual companion on the open curve]
Suppose \(w\in\mathbf Z\).  The sheaf
\[
 \mathcal F_W^c=\mathcal F_W^\vee(-w)
\]
is a type-3 datum of weight \(w\), is unramified at infinity, has no
geometric Artin--Schreier subquotient, and is arithmetically irreducible.  It
is an \(\iota\)-conjugate companion of \(\mathcal F_W\) on \(U_0\): for every
closed point \(u\), its Frobenius eigenvalue multiset becomes, under
\(\iota\), the complex-conjugate multiset of the eigenvalues of
\(\mathcal F_W\).
\end{theorem}

\begin{theorem}[Conjugate boundary characteristic polynomials]
Assume \(w\in\mathbf Z\) and put
\(\mathcal F_W^c=\mathcal F_W^\vee(-w)\).  For every
\(s\in\mathbf A^1_k\setminus U_0\),
\[
 \iota\det\!\left(1-TF_s\mid
    (\mathcal F^c_{\bar\eta_s})^{I_s}\right)
 =\overline{
 \iota\det\!\left(1-TF_s\mid
    \mathcal F_{\bar\eta_s}^{I_s}\right)}.
\]
The sheaf \(\mathcal F_W^c\) is a boundary-real companion.  The displayed
identity concerns inertia invariants and does not identify
\((V^\vee)^{I_s}\) with \((V^{I_s})^\vee\).
\end{theorem}

\begin{definition}[Paired Fourier input]
Let \(\mathcal F_W^c\) be any boundary-real companion and put
\[
 \mathcal M_W=j_{0*}(\mathcal F_W\oplus\mathcal F_W^c).
\]
This is a constructible Weil sheaf on \(A_0\), and
\(\mathcal M_W[1]\) is the direct sum of the two strict Weil middle
extensions attached to \(\mathcal F_W\) and \(\mathcal F_W^c\).
\end{definition}

\begin{theorem}[Concentration and lissity of the paired transform]
Let \(j'_0:U'_0=A'_0\setminus\{0'\}\hookrightarrow A'_0\).  For
\(\chi\in\{\psi,\psi^{-1}\}\), put
\[
 C_{\chi,W}:=j_0'^*\FT_{\chi,W}(\mathcal M_W[1])
       \in\operatorname{WeilD}^b_c(U'_0,E).
\]
Then
\[
 \mathcal H^i(C_{\chi,W})=0\quad(i\ne-1),
 \qquad
 \mathcal H^{-1}(C_{\chi,W})
 \text{ is a finite-rank lisse Weil sheaf on }U'_0.
\]
These assertions commute with the forgetful functor and with finite
extension of the coefficient field.
\end{theorem}

\begin{definition}[Paired Fourier envelope]
Define the lisse Weil sheaf \(\mathcal G'_W\) on \(U'_0\) by
\[
 \mathcal G'_W=
 \mathcal H^{-1}j_0'^*\!\left(
  \FT_{\psi,W}(\mathcal M_W[1])\oplus
  \FT_{\psi^{-1},W}(\mathcal M_W[1])
 \right).
\]
There is a unique isomorphism
\[
 \mathcal G'_W[1]\xrightarrow{\sim}j_0'^*\!\left(
  \FT_{\psi,W}(\mathcal M_W[1])\oplus
  \FT_{\psi^{-1},W}(\mathcal M_W[1])
 \right)
\]
whose degree \(-1\) cohomology map is the identity.
\end{definition}

\begin{definition}[Transformed-constituent splitting maps]
Let
\[
 u:\mathcal F_W\longrightarrow\mathcal F_W\oplus\mathcal F_W^c,
 \qquad
 p:\mathcal F_W\oplus\mathcal F_W^c\longrightarrow\mathcal F_W
\]
be the first-summand inclusion and projection, so that \(pu=1_{\mathcal
F_W}\).  Apply \(j_{0*}(-)[1]\), \(\FT_{\psi,W}\), restriction to
\(U'_0\), and degree \(-1\) cohomology; then include into, respectively
project from, the \(\psi\)-summand of \(\mathcal G'_W\).  Denote the
resulting morphisms of lisse Weil sheaves by
\[
 i_{\mathcal F'}:\mathcal F'_W\longrightarrow\mathcal G'_W,
 \qquad
 r_{\mathcal F'}:\mathcal G'_W\longrightarrow\mathcal F'_W,
\]
and put \(e_{\mathcal F'}=i_{\mathcal F'}r_{\mathcal F'}\in
\operatorname{End}(\mathcal G'_W)\).
\end{definition}

\begin{proposition}[Transformed-constituent splitting equations]
The splitting maps are Frobenius-equivariant and satisfy
\[
 r_{\mathcal F'}i_{\mathcal F'}=1_{\mathcal F'_W},\qquad
 e_{\mathcal F'}^2=e_{\mathcal F'},\qquad
 e_{\mathcal F'}i_{\mathcal F'}=i_{\mathcal F'},\qquad
 r_{\mathcal F'}e_{\mathcal F'}=r_{\mathcal F'}.
\]
The image of \(e_{\mathcal F'}\) is isomorphic to \(\mathcal F'_W\) as a
lisse Weil sheaf.  The maps and all four equations commute with finite
coefficient extension.
\end{proposition}

\begin{definition}[Real lisse Weil sheaf]
A lisse Weil \(E\)-sheaf \(\mathcal H_W\) on a scheme of finite type over
\(k\) is \emph{\(\iota\)-real} if, for every closed point \(x\),
\[
 \iota\det\!\left(1-TF_x\mid\mathcal H_{\bar x}\right)\in\mathbf R[T],
\]
where \(F_x\) is geometric Frobenius at \(x\).
\end{definition}

\begin{theorem}[Real characteristic polynomial]
For every \(n\ge1\) and every
\(x'\in U'_0(\mathbf F_{q^n})\), one has the finite Fourier trace identity
\[
 \begin{aligned}
 \Tr(F_{x',n}\mid\mathcal G'_{\bar x'})
 =-\sum_{x\in A_0(\mathbf F_{q^n})}
 &\Tr(F_{x,n}\mid\mathcal M_{\bar x})\\
 &\cdot\left[
 \psi(\Tr_{\mathbf F_{q^n}/k}(xx'))+
 \psi^{-1}(\Tr_{\mathbf F_{q^n}/k}(xx'))
 \right].
 \end{aligned}
\]
For every \(m\geq1\), one also has
\[
 \iota\Tr\!\left(F_{x',n}^{m}\mid\mathcal G'_{\bar x'}\right)
 \in\mathbf R,
\]
and
\[
 \iota\det(1-TF_{x',n}\mid\mathcal G'_{\bar x'})\in\mathbf R[T].
\]
The lisse Weil sheaf \(\mathcal G'_W\) is \(\iota\)-real, and the arithmetic
direct summand cut out by \(e_{\mathcal F'}\) is \(\mathcal F'_W\).
\end{theorem}

\begin{theorem}[Deligne's realness criterion]
Let \(\mathcal H_W\) be an \(\iota\)-real lisse Weil \(E\)-sheaf on a smooth
connected curve over \(k\), and let \(\mathcal J_W\) be a nonzero
arithmetically irreducible direct summand of \(\mathcal H_W\).  There is a
unique \(v\in\mathbf R\) such that \(\mathcal J_W\) is punctually
\(\iota\)-pure of weight \(v\).
\end{theorem}

\begin{theorem}[Purity of the transformed constituent]
Let \(\mathcal F_W\) be a type-3 datum of real weight \(w\) admitting a
boundary-real companion.  There is a unique real number \(w'\) such that
\(\mathcal F'_W\) is punctually \(\iota\)-pure of weight \(w'\).  If
\(w\in\mathbf Z\), the sheaf \(\mathcal F_W^\vee(-w)\) is a boundary-real
companion satisfying the hypothesis.
\end{theorem}

\begin{definition}[Local monodromy of the type-3 Fourier sheaf]
For the sheaf \(\mathcal F'_W\) in the type-3 Fourier-classification theorem,
its local
monodromy representation at the dual origin is
\[
  V':=\mathcal F'_{\bar\eta_{0'}}.
\]
It is a finite-dimensional \(E\)-vector space with continuous
\(W_{0'}\)-action.  Its inertia invariants are
\[
  (V')^{I_{0'}}
  :=\{v\in V':\sigma v=v\text{ for every }\sigma\in I_{0'}\}.
\]
This is the local monodromy of the whole global type-3 transform.  It is not
identified with the local Fourier transform of the boundary stalk alone.
\end{definition}

\begin{theorem}[Stationary-phase morphisms]
Let \(\mathcal F_W\) be a type-3 datum.  Global-to-local Fourier comparison
defines canonical continuous \(E\)-linear maps
\[
 a_{\mathrm{sp}}:H_c^1(A,j_*\mathcal F)\longrightarrow V',
 \qquad
 b_{\mathrm{sp}}:V'\longrightarrow V_\infty(-1).
\]
Here the first and third representations are inflated from \(W(k)\) through
\(W_{0'}\twoheadrightarrow W(k)\), using
\(W_\infty/I_\infty\simeq W(k)\) for the third.
\end{theorem}

\begin{theorem}[Inertia and Frobenius equivariance]
The maps \(a_{\mathrm{sp}}\) and \(b_{\mathrm{sp}}\) are morphisms in
\(\operatorname{Rep}_E(W_{0'})\).  Thus, for every \(g\in W_{0'}\),
\[
 \rho_{V'}(g)a_{\mathrm{sp}}=a_{\mathrm{sp}}\rho_{H_c^1}(g),
 \qquad
 \rho_{V_\infty(-1)}(g)b_{\mathrm{sp}}
       =b_{\mathrm{sp}}\rho_{V'}(g).
\]
In particular these equations hold for a geometric-Frobenius lift
\(\Phi_{0'}\), while \(I_{0'}\) acts trivially on the source of
\(a_{\mathrm{sp}}\) and the target of \(b_{\mathrm{sp}}\).
\end{theorem}

\begin{proposition}[Vanishing of the stationary-phase composite]
The composite of the two stationary-phase maps is zero:
\[
 b_{\mathrm{sp}}a_{\mathrm{sp}}=0.
\]
\end{proposition}

\begin{proposition}[Injectivity of the incoming stationary-phase map]
The map \(a_{\mathrm{sp}}\) is injective; equivalently,
\[
 \ker(a_{\mathrm{sp}})=0.
\]
\end{proposition}

\begin{proposition}[Exactness at the Fourier local monodromy]
The image of the incoming map equals the kernel of the outgoing map:
\[
 \operatorname{im}(a_{\mathrm{sp}})=\ker(b_{\mathrm{sp}})
 \quad\text{inside }V'.
\]
\end{proposition}

\begin{proposition}[Surjectivity onto the boundary module]
The map \(b_{\mathrm{sp}}\) is surjective; equivalently,
\[
 \operatorname{im}(b_{\mathrm{sp}})=V_\infty(-1).
\]
\end{proposition}

\begin{corollary}[Stationary-phase short exact sequence]
The preceding four stationary-phase statements give a short exact sequence in
\(\operatorname{Rep}_E(W_{0'})\):
\[
 0\longrightarrow H_c^1(A,j_*\mathcal F)
 \xrightarrow{a_{\mathrm{sp}}}V'
 \xrightarrow{b_{\mathrm{sp}}}V_\infty(-1)
 \longrightarrow0.
\]
\end{corollary}

\begin{theorem}[Naturality of stationary phase]
For a finite coefficient extension \(E'/E\), tensoring
\(a_{\mathrm{sp}}\) and \(b_{\mathrm{sp}}\) with \(E'\) gives the
stationary-phase maps of \(\mathcal F_W\otimes_EE'\).  For a finite ground
extension \(k_m/k\), use
\[
 \psi_m=\psi\circ\operatorname{Tr}_{k_m/k}
\]
and restrict from \(W(k)\) to
\(W(k_m)=\langle F_q^m\rangle\).  Under the Fourier base-change isomorphism,
this gives exactly the stationary-phase maps after base change to \(k_m\).
For \(a\in k^\times\), the Fourier coordinate-change isomorphism for
\(x_1=ax\) and \(x'_1=a^{-1}x'\) carries
\(a_{\mathrm{sp}},b_{\mathrm{sp}}\) to the corresponding maps in the new
dual coordinates.  The equivariance equations, zero-composite equation,
the three exactness identities, and all inverse equations are natural in
\(\mathcal F_W\) and are preserved by coefficient extension, ground-field
extension, and dual-coordinate change.
\end{theorem}

\begin{theorem}[Echelon-two local monodromy]
In the stationary-phase sequence, \(I_{0'}\) acts trivially on the two
outer terms.  Its action on \(V'\) satisfies
\[
 (\sigma-1)V'\subset a_{\mathrm{sp}}
 H_c^1(A,j_*\mathcal F),\qquad(\sigma-1)^2V'=0
 \quad(\sigma\in I_{0'}).
\]
Thus the local monodromy of \(V'\) is unipotent of echelon at most two; wild
inertia acts trivially and the action factors through the maximal pro-\(\ell\)
quotient of tame inertia.
\end{theorem}

\begin{definition}[Tame parameter at the dual origin]
Fix a compatible system of \(\ell\)-power roots of a uniformizer.  It gives
the continuous tame character
\[
 t_\ell:I_{0'}\longrightarrow\mathbf Z_\ell(1),
\]
factoring through the maximal pro-\(\ell\) tame quotient.  The chosen
geometric-Frobenius lift satisfies
\[
 t_\ell(\Phi_{0'}\sigma\Phi_{0'}^{-1})
       =q^{-1}t_\ell(\sigma).
\]
\end{definition}

\begin{theorem}[Monodromy logarithm]
There is a unique \(E\)-linear morphism
\[
 N:V'\longrightarrow V'(-1)
\]
such that, after the canonical contraction
\(E(1)\otimes_E V'(-1)\simeq V'\),
\[
 \rho_{V'}(\sigma)=1+t_\ell(\sigma)N
 \quad(\sigma\in I_{0'}).
\]
It satisfies the typed nilpotence equation
\(N(-1)\circ N=0:V'\to V'(-2)\).  If
\(\widetilde N:V'\to V'\) denotes the underlying linear operator after the
chosen trivialization of \(E(1)\), and \(F'\) is the operator induced by
\(\Phi_{0'}\), then
\[
 \widetilde N^2=0,\qquad
 F'\widetilde N(F')^{-1}=q^{-1}\widetilde N.
\]
Equivalently, \(N:V'\to V'(-1)\) is Frobenius-equivariant.
\end{theorem}

\begin{definition}[Length-two monodromy filtration]
Put
\[
 B'=\operatorname{im}(\widetilde N),\qquad
 K'=\ker(\widetilde N),\qquad Q'=V'/K'.
\]
Define the increasing filtration \(M_\bullet V'\) by
\[
 0=M_{-2}V'\subset M_{-1}V'=B'
 \subset M_0V'=K'\subset M_1V'=V'.
\]
Thus
\[
 \operatorname{Gr}_{-1}^M V'=B',\qquad
 \operatorname{Gr}_{0}^M V'=K'/B',\qquad
 \operatorname{Gr}_{1}^M V'=Q'.
\]
\end{definition}

\begin{theorem}[Graded monodromy map and inverse]
The filtration \(M_\bullet V'\) is Frobenius-stable,
\(K'=(V')^{I_{0'}}\), and \(B'\subset K'\).  The operator \(N\) induces
a Frobenius-equivariant isomorphism
\[
 \overline N:\operatorname{Gr}_{1}^M V'=Q'
       \xrightarrow{\sim}
       \operatorname{Gr}_{-1}^M V'(-1)=B'(-1).
\]
Let \(s_N:B'(-1)\to Q'\) be its inverse.  Then
\[
 s_N\overline N=1_{Q'},\qquad
 \overline N s_N=1_{B'(-1)}.
\]
Both equations commute with geometric Frobenius; in untwisted notation this
is the relation \(F'\widetilde N(F')^{-1}=q^{-1}\widetilde N\).
\end{theorem}

\begin{theorem}[Dual monodromy filtration]
On \((V')^\vee\) put
\(\widetilde N^\vee=-{}^t\!\widetilde N\).  Its monodromy filtration is
characterized by
\[
 M_i((V')^\vee)=\{\lambda\in(V')^\vee:
             \lambda(M_{-i-1}V')=0\}.
\]
Evaluation induces Frobenius-equivariant perfect pairings
\[
 \operatorname{Gr}_i^M V'\otimes_E
 \operatorname{Gr}_{-i}^M (V')^\vee\longrightarrow E
 \qquad(i=-1,0,1).
\]
Under these pairings the transpose of \(\overline N\) is the graded map
induced by \(\widetilde N^\vee\), and the transpose of \(s_N\) is its
inverse.  For the local Verdier dual the same statements hold after the
common Tate twist, with target \(E(1)\).
\end{theorem}

\begin{theorem}[Factorization through inertia invariants]
The map \(a_{\mathrm{sp}}\) has image in
\(K'=(V')^{I_{0'}}\).  Its corestriction is the Frobenius-equivariant map
\[
 \overline a_{\mathrm{sp}}:H_c^1(A,j_*\mathcal F)
       \longrightarrow(V')^{I_{0'}}.
\]
\end{theorem}

\begin{theorem}[Inverse on inertia invariants]
There is a canonical Frobenius-equivariant stationary-phase map
\[
 c_{\mathrm{inv}}:(V')^{I_{0'}}\longrightarrow
       H_c^1(A,j_*\mathcal F)
\]
satisfying the two inverse equations
\[
 c_{\mathrm{inv}}\overline a_{\mathrm{sp}}
       =1_{H_c^1(A,j_*\mathcal F)},\qquad
 \overline a_{\mathrm{sp}}c_{\mathrm{inv}}
       =1_{(V')^{I_{0'}}}.
\]
The map and both equations commute with finite coefficient extension and
finite ground-field extension.
\end{theorem}

\begin{theorem}[Compactly supported cohomology and inertia invariants]
The maps \(\overline a_{\mathrm{sp}}\) and \(c_{\mathrm{inv}}\) give a
Frobenius-equivariant isomorphism
\[
 H_c^1(A,j_*\mathcal F)\xrightarrow{\sim}(V')^{I_{0'}}.
\]
The map induced by \(b_{\mathrm{sp}}\) on the quotient is the
Frobenius-equivariant isomorphism
\[
 \overline b_{\mathrm{sp}}:Q'=V'/(V')^{I_{0'}}
       \xrightarrow{\sim}V_\infty(-1).
\]
Its inverse \(s_{\mathrm{sp}}:V_\infty(-1)\to Q'\) satisfies
\[
 s_{\mathrm{sp}}\overline b_{\mathrm{sp}}=1_{Q'},\qquad
 \overline b_{\mathrm{sp}}s_{\mathrm{sp}}=1_{V_\infty(-1)}.
\]
In particular, \(Q'\ne0\).
\end{theorem}

\begin{theorem}[Upper bound on boundary invariants]
Assume that the transformed sheaf \(\mathcal F'_W\) is punctually
\(\iota\)-pure of weight \(w'\).  The Frobenius-stable subspace
\((V')^{I_{0'}}\) has \(\iota\)-weights at most \(w'\).
\end{theorem}

\begin{theorem}[Upper bound on the monodromy quotient]
Under the same purity hypothesis, and for the echelon-two local monodromy
of the type-3 Fourier transform, the quotient
\[
 Q'=V'/(V')^{I_{0'}}
\]
has \(\iota\)-weights at most \(w'+1\).  Equivalently, each Frobenius
eigenvalue \(\beta\) on \(Q'\) satisfies
\[
 |\iota(\beta)|\le q^{(w'+1)/2}.
\]
\end{theorem}

\begin{theorem}[Contragredient bound for the dual quotient]
Under the same purity hypothesis, the dual monodromy filtration of the
ordinary contragredient representation \((V')^\vee\) has a degree \(-1\)
graded piece canonically isomorphic to \((Q')^\vee\).  This Frobenius module has
\(\iota\)-weights at most \(-w'-1\).  The isomorphism is compatible with the transpose of
\(\overline N\), the transpose inverse of \(s_N\), and finite coefficient
extension.
\end{theorem}

\begin{theorem}[Lower bound on the monodromy quotient]
The quotient \(Q'\) has \(\iota\)-weights at least \(w'+1\).  Every
eigenvalue \(\beta\) on \(Q'\) satisfies
\[
 |\iota(\beta)|\ge q^{(w'+1)/2}.
\]
\end{theorem}

\begin{corollary}[Purity of the monodromy quotient]
The nonzero quotient \(Q'\) is pure of weight \(w'+1\).
\end{corollary}

\begin{theorem}[Fourier weight shift]
Let \(\mathcal F_W\) have weight \(w\), and suppose the transformed sheaf
has weight \(w'\).  The stationary-phase isomorphism is an isomorphism of
nonzero Frobenius modules
\[
 \overline b_{\mathrm{sp}}:Q'\xrightarrow{\sim}V_\infty(-1).
\]
The source is pure of weight \(w'+1\), the target is pure of weight \(w+2\),
and
\[
 w'+1=w+2,
 \qquad w'=w+1.
\]
\end{theorem}

\begin{theorem}[Upper weight bound on compactly supported cohomology]
Under the hypotheses of the Fourier weight-shift theorem,
\[
 (V')^{I_{0'}}\quad\text{has weights at most }w'=w+1.
\]
Through the Frobenius-equivariant stationary-phase isomorphism, the same
bound holds for
\[
 H_c^1(A,j_*\mathcal F)\xrightarrow{\sim}(V')^{I_{0'}}.
\]
\end{theorem}

\section{Laumon's affine-line estimate}

\begin{theorem}[Affine-line upper-weight theorem]
Let \(w\in\mathbf R\), let \(j_0:U_0\hookrightarrow\Aone_k\) be a dense
open immersion, and let \(\mathcal F_W\) be a type-3 datum of weight \(w\)
with coefficients in \(E\).  Assume either
\begin{enumerate}
\item the type-3 Fourier sheaf \(\mathcal F'_W\) is independently known to
      be punctually \(\iota\)-pure of some real weight \(w'\); or
\item \(\mathcal F_W\) admits a boundary-real \(\iota\)-companion on
      \(\mathbf A^1\).
\end{enumerate}
In the second case \(\mathcal F'_W\) is punctually \(\iota\)-pure.  In
either case its weight is \(w'=w+1\), and
\[
  H_c^i(A,j_*\mathcal F)=0\qquad\text{for every }i\ne1,
\]
the \(E\)-vector space \(H_c^1(A,j_*\mathcal F)\) is finite-dimensional,
and every eigenvalue \(\alpha\in\overline E\) of geometric Frobenius on it
satisfies
\[
  |\iota(\alpha)|\le q^{(w+1)/2}.
\]
Equivalently, \(H_c^1(A,j_*\mathcal F)\) has \(\iota\)-weights at most
\(w+1\).  The vanishing and finite-dimensionality assertions are separate
parts of the theorem and are not consequences of the numerical inequality.
\end{theorem}

\begin{corollary}[Integral-weight affine-line theorem]
If \(w\in\mathbf Z\), the sheaf \(\mathcal F_W^\vee(-w)\) is a boundary-real
companion.  Every type-3 datum of integral weight satisfies the conclusions
of the affine-line upper-weight theorem without an additional companion
hypothesis.
\end{corollary}

\begin{corollary}[All-embeddings form]
Suppose that the irreducibility, unramifiedness, and geometric
nonconstancy conditions in the definition of a type-3 datum hold and that
\(\mathcal F_W\) is punctually pure of the same real
weight \(w\) for every embedding
\(\iota:\overline E\hookrightarrow\mathbf C\).  Assume for each \(\iota\)
either that the transformed sheaf is independently punctually
\(\iota\)-pure, or that a boundary-real \(\iota\)-companion on
\(\mathbf A^1\) exists.  When \(w\in\mathbf Z\), the second alternative
holds with \(\mathcal F_W^\vee(-w)\).  Then every Frobenius
eigenvalue \(\alpha\) on \(H_c^1(A,j_*\mathcal F)\) satisfies
\[
 |\iota(\alpha)|\le q^{(w+1)/2}
 \quad\text{for every }\iota:\overline E\hookrightarrow\mathbf C.
\]
Thus this cohomology group has the all-embeddings upper weight bound
\(w+1\).  A bound for one chosen embedding does not give this conclusion.
\end{corollary}

\begin{proposition}[Coefficient-field extension and descent]
For every finite extension \(E'/E\), there are canonical
Frobenius-equivariant isomorphisms
\[
 H_c^i(A,j_*\mathcal F)\otimes_EE'
 \xrightarrow{\sim}
 H_c^i\bigl(A,j_*(\mathcal F\otimes_{E_U}E'_U)\bigr).
\]
Ordinary middle extension, Fourier transform, and the local exact sequence
commute with this scalar extension.  The characteristic polynomial of
Frobenius after extension is the image in \(E'[T]\) of the original
polynomial, so an upper \(\iota\)-weight bound holds before scalar extension
if and only if it holds after scalar extension for the induced embedding.

For a \(\Qlb\)-coefficient Karoubi triple \((E,K,e)\), compactly supported
cohomology is the image of the Frobenius-equivariant idempotent
\(H_c^i(e)\) on the finite-stage group \(H_c^i(A,K)\).  Hence the upper
weight bound passes from the finite-stage group to the direct summand and
is independent of the chosen common finite coefficient stage.
\end{proposition}

\begin{proposition}[Ground-field extension and descent]
After replacing \(k=\Fq\) by \(k_n=\mathbf F_{q^n}\), the underlying
geometric cohomology spaces and the local exact sequence are unchanged.  The
Fourier kernel is formed with
\[
 \psi_n=\psi\circ\operatorname{Tr}_{k_n/k},
\]
and the natural Fourier base-change isomorphism identifies every transform,
splitting map, stationary-phase map, and inverse map with its restriction to
\(W(k_n)=\langle F_q^n\rangle\).  Geometric Frobenius is therefore replaced
by its \(n\)-th power.  Consequently
\[
  |\iota(\alpha)|\le q^{v/2}
  \quad\Longleftrightarrow\quad
  |\iota(\alpha^n)|\le(q^n)^{v/2}.
\]
The affine-line estimate is therefore invariant under finite extension of
the ground field.
\end{proposition}

\section{References}

The statement and normalization follow Deligne's theory of weights and
Laumon's Fourier transform.  The pro-etale coefficient language may be
compared with the ordinary etale theory by the constructibility and
comparison theorems of Bhatt--Scholze; the two realizations agree for the
finite-stage lisse objects used here.

\begin{thebibliography}{9}
\bibitem{DeligneWeilII}
P. Deligne, \emph{La conjecture de Weil II},
Publications Mathematiques de l'IHES \textbf{52} (1980), 137--252.

\bibitem{LaumonFourier}
G. Laumon, \emph{Transformation de Fourier, constantes d'equations
fonctionnelles et conjecture de Weil},
Publications Mathematiques de l'IHES \textbf{65} (1987), 131--210.

\bibitem{BhattScholze}
B. Bhatt and P. Scholze, \emph{The pro-etale topology for schemes},
Asterisque \textbf{369} (2015), 99--201.
\end{thebibliography}

\end{document}
